% SPDX-License-Identifier: CC-BY-4.0
\documentclass[11pt,letterpaper]{article}

\usepackage[T1]{fontenc}
\usepackage{amsmath,amssymb}
\usepackage{tikz}
\usepackage{pgfplots}
\pgfplotsset{compat=1.18}
\usepackage{algorithm}
\usepackage{algpseudocode}
\usepackage{booktabs}
\usepackage{enumitem}
\usepackage[hidelinks]{hyperref}

\newcommand{\N}{\mathbb{N}}
\newcommand{\Z}{\mathbb{Z}}
\newcommand{\Q}{\mathbb{Q}}
\newcommand{\R}{\mathbb{R}}
\newcommand{\C}{\mathbb{C}}
\newcommand{\Map}[3]{#1 \colon #2 \to #3}
\newcommand{\Mod}[1]{#1\text{-}\mathrm{Mod}}
\newcommand{\GroupHom}{\mathrm{Hom}_{\mathbf{Grp}}}
\newcommand{\RingHom}{\mathrm{Hom}_{\mathbf{Ring}}}

\newtheorem{theorem}{Theorem}
\newtheorem{lemma}[theorem]{Lemma}
\newtheorem{definition}[theorem]{Definition}
\newtheorem{remark}[theorem]{Remark}
\newenvironment{proof}{\noindent\textit{Proof.}}{\hfill\(\square\)}

\title{The positive additive map}
\author{}
\date{\today}

\begin{document}
\maketitle

\begin{abstract}
For positive additive generalized Collatz maps, the greatest common divisor
of the offset and divisor determines access to division. We prove the unique
access index, recover its entire raw block, and construct an explicit finite
first-division core. A signed reference subtraction gives an exact geometric
entry budget, including its conversion to the original instruction count.
Inaccessible states have explicit affine-only escaping tails.
\end{abstract}

\section{The additive map}

% The positive domain is retained throughout the first-division construction.
\begin{definition}[Additive dynamics]
Fix integers \(c\geq1\) and \(D\geq2\). On positive integers define
\[
 T(n)=\begin{cases}n/D,&D\mid n,\\ n+c,&D\nmid n.\end{cases}
\]
Write \(T^j\) for the \(j\)-fold iterate, with \(T^0(n)=n\). Put
\[
 g=\gcd(c,D),\qquad L=D/g,\qquad
 B=c(L-1),\qquad N=\left\lfloor\frac{B}{D-1}\right\rfloor.
\]
A state is accessible when \(g\mid n\), and inaccessible otherwise.
\end{definition}

\section{Access and raw reconstruction}

% Cancellation after dividing by the gcd supplies existence and uniqueness together.
\begin{theorem}[Unique access]
For a positive integer \(n\), an integer \(j\geq0\) with \(D\mid n+jc\)
exists if and only if \(g\mid n\). When \(g\mid n\), there is exactly one
\(0\leq\ell<L\) with this property. No \(0\leq j<\ell\) satisfies
\(D\mid n+jc\).
\end{theorem}

\begin{proof}
If \(D\mid n+jc\), then \(g\) divides both \(n+jc\) and \(jc\), so it
divides \(n\). Conversely, suppose \(n=ga\), and write \(c=gb\), \(D=gL\).
Then \(\gcd(b,L)=1\). The congruence \(a+jb\equiv0\pmod L\) therefore
has exactly one residue class of solutions modulo \(L\), obtained by
multiplying \(-a\) by the inverse of \(b\). Choose its representative
\(\ell\) in \(0,\ldots,L-1\). If \(L=1\), that set consists only of
\(0\) and the congruence is automatic. If another \(j\) in this interval
were a solution, cancellation of the unit \(b\) would give
\(j\equiv\ell\pmod L\), hence \(j=\ell\). An earlier solution would
be another representative in the same interval, which is impossible.
\end{proof}

% The guards establish the exact raw clock rather than just the endpoint equation.
\begin{theorem}[Exact first-division block]
For an accessible state \(n\), let \(\ell\) be the index above and set
\[
 H(n)=\frac{n+\ell c}{D}.
\]
Then \(H(n)\) is a positive integer and
\[
 T^j(n)=n+jc\quad(0\leq j\leq\ell),\qquad
 T^{\ell+1}(n)=H(n).
\]
For an inaccessible state, \(T^j(n)=n+jc\) for every \(j\geq0\), and this
full sequence tends to infinity.
\end{theorem}

\begin{proof}
Every pre-division state \(n+jc\), \(j<\ell\), fails the divisibility guard,
so induction gives the stated affine prefix. At index \(\ell\) the division
guard holds and the next raw instruction divides by \(D\). The numerator
is a positive multiple of \(D\), proving positivity of the endpoint.
For an inaccessible state, the preceding theorem rules out division at every
affine index. Induction therefore gives the infinite affine tail. Since
\(c\geq1\), it is at least \(n+j\) and tends to infinity.
\end{proof}

\section{Finite first-division cores}

% Strict integer descent removes all dependence on arbitrarily large starting states.
\begin{theorem}[Finite core]
Every accessible state satisfies \(D H(n)\leq n+B\). If \(n>N\), then
\(H(n)<n\); if \(n\leq N\), then \(H(n)\leq N\). Also \(N\leq c\).
Every positive seed reaches a positive state at most \(N\), or reaches an
inaccessible state, at a finite raw time.
\end{theorem}

\begin{proof}
The bound \(\ell\leq L-1\) gives
\(D H(n)=n+\ell c\leq n+B\). If \(n>N\), integrality gives
\((D-1)n>B\), whence \(n+B<Dn\) and \(H(n)<n\). If \(n\leq N\), use
\(B<(D-1)(N+1)\) to obtain
\[
 D H(n)\leq N+B<D(N+1).
\]
Thus the integer \(H(n)\) is at most \(N\). Since \(L\leq D\), we have
\(B\leq c(D-1)\), giving \(N\leq c\).

For the finite-time assertion, use strong induction on the positive seed.
A seed already in the core or inaccessible needs no instruction. Otherwise
it is accessible and above \(N\), so its first-division endpoint is a
smaller positive integer. Apply the induction hypothesis to that endpoint
and prepend its \(\ell+1\) raw instructions. This constructs a finite raw
witness with the same core-or-inaccessible conclusion. The induction uses
only smaller positive integers. Core and inaccessible states provide its
base cases.
\end{proof}

% Coprimality removes traps; an empty core instead forces an affine escaping tail.
\begin{theorem}[Boundary cases]
If \(g=1\), every positive seed reaches the set \(\{1,\ldots,c\}\).
All its accessible first-division edges stay within this set. If \(N=0\),
every positive raw orbit tends to infinity.
\end{theorem}

\begin{proof}
For \(g=1\), every state is accessible, \(L=D\), and \(N=c\).
The finite-core theorem therefore cannot terminate at an inaccessible state
and gives the claimed entry and invariance. If \(N=0\), no positive state
is in the core. The same theorem supplies a raw time \(t\) at which the
state \(x=T^t(n)\) is inaccessible. Its subsequent orbit is \(x+jc\).
For every sufficiently large raw time \(k\), write \(k=t+j\); then
\(T^k(n)=x+(k-t)c\), proving full-sequence divergence, not merely an
unbounded subsequence.
\end{proof}

\section{A geometric instruction budget}

% Reference subtraction cancels the repeated additive displacement without an error term.
\begin{definition}[Shifted coordinate]
For any integer state set
\[
 \delta(n)=(D-1)n-B,\qquad \varepsilon=\delta(N+1).
\]
Writing \(B=(D-1)N+r\), with \(0\leq r<D-1\), gives
\(\varepsilon=D-1-r\), hence \(1\leq\varepsilon\leq D-1\).
\end{definition}

% The same integer threshold works for every legal path, including paths ending at a trap.
\begin{theorem}[Geometric entry and raw recovery]
Let \(h\geq0\) satisfy
\[
 D^h\varepsilon>\delta(n).
\]
Then the raw orbit of \(n\) reaches the core or an inaccessible state at some
time \(t\leq hL\). Such an \(h\) always exists; the least one is an exact
integer geometric budget.
\end{theorem}

\begin{proof}
For an accessible edge \(n\mapsto m=H(n)\), multiply the endpoint inequality
by \(D-1>0\) and subtract \(DB\):
\[
 D\delta(m)\leq(D-1)(n+B)-DB=\delta(n).
\]
For any path \(n_0=n,n_1,\ldots,n_h\) of legal first-division edges,
induction, multiplying each inequality by the positive power \(D^i\), gives
\(D^h\delta(n_h)\leq\delta(n)\). If \(n_h>N\), then
\(\delta(n_h)\geq\delta(N+1)=\varepsilon\), contradicting the assumed
strict budget inequality. Hence the last endpoint lies in the core.

Construct edges successively only while the current state is accessible and
outside the core. If construction stops before \(h\) edges, its endpoint
already has the required alternative. Otherwise the preceding inequality
forces the \(h\)-th endpoint into the core. This reasoning also covers
\(h=0\): the budget inequality then directly gives \(n\leq N\).

If \(k\leq h\) edges were used, let their access indices be
\(\ell_0,\ldots,\ell_{k-1}\). Their exact raw expansion has length
\[
 t=\sum_{i=0}^{k-1}(\ell_i+1)\leq kL\leq hL.
\]
Composition of the exact raw-block identities proves that its endpoint is
\(T^t(n)\). No guard or intermediate instruction is discarded.
Finally, \(\varepsilon\geq1\), \(D\geq2\), and the integer powers of
\(D\) are unbounded; hence a suitable finite \(h\) exists. The least one
is obtained by repeated exact multiplication and comparison.
\end{proof}

\section{Finite graphs and examples}

The core is a finite set of first-division states, not a raw-state bound.
Each accessible core state has its unique edge \(n\mapsto H(n)\); an
inaccessible state retains its affine-only tail instead of being given a
spurious fixed-point edge. The exact raw word and its length accompany
every first-division edge. This construction does not assert a classification
of the resulting cycles.

For \((c,D)=(6,4)\), the constants are \(g=2\), \(L=2\), \(N=2\).
The accessible state \(2\) has raw cycle \(2\mapsto8\mapsto2\), while
\(4\mapsto1\) reaches the inaccessible tail \(1,7,13,\ldots\).
For \((c,D)=(2,4)\), the radius is zero, so every positive seed escapes.
With \(c=1\), \(D=37\), the seed \(53\cdot37^{300}\) reaches \(1\)
in \(302\) first-division edges and \(358\) raw instructions: after the
first \(300\) divisions, \(53\) needs \(21\) additions to reach \(74\),
then division gives \(2\), which needs \(35\) additions before division
gives \(1\).

\section{Formal and computational statements}

The Lean development proves access and uniqueness, all raw guards and block
identities, strict descent, core invariance, finite raw recovery, the shifted
path inequality, the unconditional geometric raw-time budget, coprime core
entry, and full-sequence escape for the zero-core case.
The arithmetic bounds also permit \(c=0\); the divergence statements
explicitly require \(c>0\). The present positive-offset theorem makes no
negative-divisor or signed-state assertion.

The Go executable uses exact integer modular inversion and arbitrary-precision
values. A certificate records its seed, first core-or-inaccessible endpoint,
raw instruction count, edge count, least geometric budget, and alternative.
Verification recomputes these fields without changing the supplied
certificate. The bounded family comprises \(32\cdot47\cdot128=192{,}512\)
certificates, with the large entry above as a separate check. These finite
checks concern the implementation; the universal statements follow from the
proofs, independently of the tested range.

\end{document}
